
% ============================================================
% EVALUACIÓN
%
% Esta sección presenta las actividades evaluativas y sus
% porcentajes mediante una tabla de dos columnas.
% ============================================================

% Crea el encabezado numerado de la sección.
\section{Evaluación}

La evaluación del curso se distribuirá de la siguiente manera:

% tabularx crea una tabla que ocupa todo el ancho del texto.
%
% La primera columna X:
% - utiliza el espacio disponible;
% - alinea el contenido a la izquierda.
%
% La segunda columna:
% - tiene un ancho fijo;
% - alinea los porcentajes a la derecha.
\begin{tabularx}{\textwidth}{@{}
    >{\raggedright\arraybackslash}X
    >{\raggedleft\arraybackslash}p{0.18\textwidth}
@{}}

% \toprule crea la línea superior.
\toprule

% Encabezados de las columnas.
\textbf{Actividad} & \textbf{Porcentaje} \\

% \midrule separa el encabezado del contenido.
\midrule

% \, añade un espacio pequeño antes del signo de porcentaje.
% \% permite imprimir el carácter especial «%».
Actividad evaluativa 1 & XX\,\% \\

Actividad evaluativa 2 & XX\,\% \\

Actividad evaluativa 3 & XX\,\% \\

% \bottomrule crea la línea inferior.
\bottomrule

% Finaliza la tabla.
\end{tabularx}

% Sustituye XX por valores numéricos y comprueba que el total
% corresponda al 100 %.
Los porcentajes deben revisarse y completar el 100\,\%.